\scanpage{145}
\deleted{\unsure{Dans le domaine de la} \illegible{1}}

Soit $M\in{\unsure{\mathcal A|A}}^G$, alors pour $\forall$ entier $n$, $\bigotimes^n M$ est un
élément de $\mathcal A^{G\times\unsure{\mathfrak S|\mathcal S}_n}=(\mathcal A^G)^{\mathfrak S_n}$. En \deleted{\unsure{tout}} \unsure{plus bref|par abr.},
on notera cet élément par $T^n(M)$. Ainsi $T^n$ est
un foncteur (non plus additif !)
% slanted
\margin{Expliquer la \unsure{variation|opérations} : on considère $\boxtimes^n M\in\mathcal A^{G^n}$, où $\mathfrak S_n$ opère dans $G^n$ et dans $\boxtimes^n M$\ldots}
% slanted
\margin{Dire que si un \unsure{gpe|groupe} fini $G'$ opère dans un élément de $\mathcal A^G$, on obtient un élément de $\mathcal A^{\unsure{G\rtimes G'|G\times G'}}$}
\[
T^n:\mathcal A^G\longrightarrow(\mathcal A^G)^{\mathfrak S_n} \tag*{(15)}
\]
Sa propriété essentielle est
\[
\boxed{T^n(M+M')\simeq\sum_{p+q=n}T^p(M)*T^q(M')} \tag*{(16)}
\]
% slanted
\margin{et propriété analogue pour les \uline{\unsure{diviseurs|divisés}} ?}
De plus, $T^n$ commute aux $\mathcal A^u$ : Soit $u:G\longrightarrow G'$,
alors le diagramme de foncteurs
\diagram{\begin{tikzcd}[ampersand replacement=\&,column sep=huge]
\text{(17)} \& \mathcal A^{G'} \arrow[r,"T^n_{G'}"] \arrow[d,"\mathcal A^u"'] \& (\mathcal A^{G'})^{\mathfrak S_n} \arrow[d,"\mathcal A^{u\times 1_{\mathfrak S_n}}"]\\
\&\mathcal A^G \arrow[r,"T^n_G"] \&(\mathcal A^G)^{\mathfrak S_n}
\end{tikzcd}}
est commutatif (à préciser les compatibilités
\deleted{\unsure{des}} éventuelles d'hom. fonctoriels).

Considérons
\[
v_{p,q}:\mathfrak S_p\times\mathfrak S_q\longrightarrow\mathfrak S_{pq} \tag*{(18)}
\]
défini par $v_{p,q}(\sigma,\tau)(i)=\alpha_{p,q}(\sigma)\,u_{\underbrace{\tau,\ldots,\tau}_{p}}\,\alpha_{p,q}(\sigma)^{-1}$
où $\alpha_{p,q}\deleted{(\sigma)}:\mathfrak S_p\longrightarrow\mathfrak S_{pq}$ est défini par
\[
\left\{\begin{aligned}
&\alpha_{p,q}(\sigma)\text{ est croissante dans les intervalles } ]iq,i(q+1)]\\
&\text{et l'applique sur l'intervalle}\\
&]\sigma(i)q,\sigma(i)(q+1)]
\end{aligned}\right. \tag*{(19)}
\]

Sous ces conditions, on a
\[
T^m(T^n(M))\simeq\deleted{\unsure{v_{m,n}}}\,\mathcal A^{v_{m,n}}(T^{mn}(M)) \tag*{(20)}
\]
(les deux mbres dans $\mathcal A^{G\times\mathfrak S_m\times\mathfrak S_n}$)
% archive sheet number: 144
% printed form (vertical paper watermark)
\unsure{NTINE} \tikz[baseline=-.2ex,scale=.17]{\draw (0,0)--(1.6,0)--(1.8,.33)--(.5,.33)--(.32,.13)--(.92,.13);} TYPEWRITER \tikz[baseline=-.4ex,scale=.15]{\foreach \a in {0,45,...,315}{\draw (\a:.52) circle (.22);} \draw (0,0) circle (.21);} \unsure{ARGENTINE|ARGENTIN}
